\FloatBarrier
\section{Singlet--Triplet Qubit Control}
\label{app:st-control}

A singlet--triplet qubit stores its logical state in the singlet and
unpolarized triplet of two electrons. Exchange separates their energies,
while a Zeeman-energy difference mixes them. Two such qubits can also
acquire a conditional phase through the Coulomb interaction between
their state-dependent charge distributions. This appendix develops the
one-qubit rotations and this capacitive two-qubit gate, including the
residual interaction during nominally local operations.

Logical-state preparation and spin-to-charge readout are described in
Appendices~\ref{app:state-preparation} and~\ref{app:spin-readout}.

\subsection{Single-qubit gates for singlet--triplet qubits}
\label{sec:calibration-st-rotations}
\label{sec:measurement-st-exchange}

The logical Bloch sphere has two control mechanisms. A difference
between the two local Zeeman splittings mixes $S$ and $T_0$ and
therefore rotates about the logical $x$ axis. Exchange separates
their energies and rotates about the logical $z$ axis. The gradient
usually remains present when exchange is turned on, so an exchange
pulse generally rotates about a tilted axis rather than exactly
about $z$. We will measure the two terms, construct $x$ and $z$
rotations from them, and then combine those rotations into an
arbitrary one-qubit gate.

\subsubsection{Calibrating the two terms in the logical Hamiltonian}

Use $\ket{0_{\mathrm{ST}}}=\ket S$ and
$\ket{1_{\mathrm{ST}}}=\ket{T_0}$, and define
\[
 Z_{\mathrm{ST}}=\ket S\!\bra S-\ket{T_0}\!\bra{T_0},
 \qquad
 X_{\mathrm{ST}}=\ket S\!\bra{T_0}+\ket{T_0}\!\bra S.
\]
At a fixed voltage setting, the interdot barrier determines the
exchange splitting $J$, while the difference
$\Delta E_Z=E_{Z,1}-E_{Z,2}$ between the local Zeeman splittings
gives the transverse term. The logical Hamiltonian is
\begin{equation}
 H_{\mathrm{ST}}
 =-\frac{J}{2}Z_{\mathrm{ST}}
  +\frac{\Delta E_Z}{2}X_{\mathrm{ST}}.
 \label{eq:calibration-st-axis}
\end{equation}
Holding this setting for a time $\tau$ makes the two coefficients
constant. With $E=\sqrt{J^2+\Delta E_Z^2}$, the corresponding
propagator is
\begin{equation}
 U_{\mathrm{ST}}(\tau)
 =\cos\frac{E\tau}{2\hbar}I
  -i\sin\frac{E\tau}{2\hbar}
       \frac{\Delta E_ZX_{\mathrm{ST}}-JZ_{\mathrm{ST}}}{E}.
 \label{eq:calibration-st-actual-pulse}
\end{equation}
Thus one voltage setting specifies the rotation axis
$(\Delta E_Z,0,-J)$, and the hold time specifies the rotation
angle $E\tau/\hbar$. Gate construction requires both coefficients
to be measured.

To measure $|\Delta E_Z|$, prepare the separated singlet using
Section~\ref{sec:calibration-st-initialization}, set $J$ negligible,
and hold for a variable time $\tau$. Equation~\eqref{eq:calibration-st-actual-pulse}
then gives
\begin{equation}
 P_{T_0|S}(\tau)
 =\sin^2\!\left(\frac{\Delta E_Z\tau}{2\hbar}\right).
 \label{eq:calibration-st-gradient-oscillation}
\end{equation}
Use the pair readout of Section~\ref{sec:calibration-st-readout}
after every hold. If the measured oscillation period is $T_\Delta$,
then $|\Delta E_Z|=h/T_\Delta$.

Measuring $J$ requires an input with both $S$ and $T_0$ components,
because a pure singlet is an eigenstate of the exchange term. At
$J\simeq0$, use the calibrated gradient to apply
$R_x(\pi/2)$ to $\ket S$. Pulse to the voltage at which $J$ is to
be measured, hold for a variable time $\tau$, return to the
gradient setting, and apply $R_x(-\pi/2)$ before pair readout. If
the hardware cannot reverse the gradient, the inverse is the
equivalent positive rotation $R_x(3\pi/2)$. When the gradient is
negligible during the exchange hold, the singlet-return probability
is
\begin{equation}
 P_S(\tau)=\frac{1+\cos\theta_J}{2},
 \qquad
 \theta_J=\frac1\hbar\int J(t)\,dt.
 \label{eq:measurement-st-exchange-interference}
\end{equation}
For a constant plateau, the oscillation period is $T_J=h/J$.
Repeating the measurement at different compensated barrier voltages
determines $J(V)$. If the gradient is appreciable during the hold,
fit the measured return probability to
\[
 \left|\bra S R_x(-\pi/2)
 U_{\mathrm{ST}}(\tau)R_x(\pi/2)\ket S\right|^2
\]
using the independently measured $\Delta E_Z$; the fitted parameter
is then $J$. Propagators for the rising and falling edges can be
included in the same product.

\subsubsection{Constructing logical \texorpdfstring{$X$}{X} and
\texorpdfstring{$Z$}{Z} rotations}

With $J\simeq0$, Eq.~\eqref{eq:calibration-st-axis} reduces to
$H_{\mathrm{ST}}=\Delta E_ZX_{\mathrm{ST}}/2$. A hold of duration
$\tau_x$ therefore implements
\begin{equation}
 R_x(\theta_x)
 =\exp\!\left(-\frac{i\theta_x}{2}X_{\mathrm{ST}}\right),
 \qquad
 \theta_x=\frac{\Delta E_Z\tau_x}{\hbar}.
 \label{eq:calibration-st-x-rotation}
\end{equation}
This supplies the first logical rotation axis.

For a target $R_z(-\theta_z)$, an exchange-dominated pulse suggests
the duration $\tau_z=\hbar\theta_z/J$. The pulse is only a pure
$z$ rotation when the gradient term can be neglected. Starting from
$S$, the unwanted $S\rightarrow T_0$ probability during the actual
tilted rotation is
\begin{equation}
 P_{T_0|S}(\tau_z)=
 \frac{\Delta E_Z^2}{J^2+\Delta E_Z^2}
 \sin^2\!\left(\frac{E\tau_z}{2\hbar}\right).
 \label{eq:calibration-st-axis-leakage}
\end{equation}
This is a logical rotation error, not population outside the
$\{S,T_0\}$ space. Requiring it to remain below $q$ for every
duration gives the sufficient condition
\[
 J\geq |\Delta E_Z|\sqrt{q^{-1}-1}.
\]
Even when this bound is satisfied, the actual rotation angle is
$E\tau_z/\hbar$, so the duration is refined by measuring the phase
of a superposition rather than by using $J\tau_z/\hbar$ alone.

If the accessible $J$ does not satisfy this condition, a longer
pulse does not help: it increases the angle but leaves the axis
tilted. Rotations about two accessible axes can instead synthesize
the required $z$ rotation. For a concrete construction, take
$\Delta E_Z=\Delta>0$ and suppose the controls can set $J=0$ or
$J=\Delta$. The first setting gives $R_x(\theta)$ in time
$\hbar\theta/\Delta$. The second setting has the axis
$\mathbf n=(1,0,-1)/\sqrt2$ and gives
$A=R_{\mathbf n}(\pi)$ in time
$\pi\hbar/(\sqrt2\Delta)$. Since this $\pi$ rotation maps
$X_{\mathrm{ST}}$ to $-Z_{\mathrm{ST}}$,
\begin{equation}
 A R_x(\theta)A=-R_z(-\theta).
 \label{eq:calibration-st-composite-z}
\end{equation}
The minus sign is a global phase. The three voltage holds
$J=\Delta$, $J=0$, and $J=\Delta$, with the durations just given,
therefore implement the target $z$ rotation in this ideal model.
For finite switching edges or voltage-dependent $\Delta E_Z$, use
the measured Hamiltonian throughout the waveform and multiply the
corresponding propagators before adjusting the three durations.

\subsubsection{Constructing an arbitrary logical one-qubit gate}

The calibrated $x$ and $z$ rotations complete the single-qubit
control problem. Up to a global phase, any target unitary on the
$\{S,T_0\}$ space can be written as
\begin{equation}
 U=R_z(\alpha)R_x(\beta)R_z(\gamma).
 \label{eq:calibration-st-euler-decomposition}
\end{equation}
The rightmost operation acts first. Implement $R_x(\beta)$ at the
$J\simeq0$ setting using Eq.~\eqref{eq:calibration-st-x-rotation},
and implement each $z$ rotation either with a calibrated
exchange-dominated hold or with the composite construction in
Eq.~\eqref{eq:calibration-st-composite-z}. This sequence turns the
measured functions $J(V)$ and $\Delta E_Z(V)$ into the voltage
settings and durations of a specified logical gate.

\subsection{Two-qubit gates for singlet--triplet qubits}
\label{sec:calibration-controlled-phase}

Suppose a pulse does not transfer population between the logical
basis states $\ket{ab}=\ket a_A\ket b_B$. In the ordered basis
$\{\ket{00},\ket{01},\ket{10},\ket{11}\}$, its unitary is
diagonal. Direct factorization gives
\begin{equation}
 \begin{aligned}
 &\operatorname{diag}\!\left(
   e^{i\chi_{00}},e^{i\chi_{01}},
   e^{i\chi_{10}},e^{i\chi_{11}}\right)\\
 &\quad=e^{i\chi_{00}}
 \left[
  \operatorname{diag}\!\left(1,e^{i(\chi_{10}-\chi_{00})}\right)
  \otimes
  \operatorname{diag}\!\left(1,e^{i(\chi_{01}-\chi_{00})}\right)
 \right]\\
 &\qquad\times
 \operatorname{diag}\!\left(
   1,1,1,e^{i(\chi_{11}-\chi_{10}-\chi_{01}+\chi_{00})}
  \right).
 \end{aligned}
 \label{eq:calibration-phase-factorization}
\end{equation}
Denote the phases in the first local factor, the second local
factor and the final factor by
\[
 \alpha=\chi_{10}-\chi_{00},\qquad
 \beta=\chi_{01}-\chi_{00},\qquad
 \Theta=\chi_{11}-\chi_{10}-\chi_{01}+\chi_{00}.
\]

The factor $e^{i\chi_{00}}$ is a global phase. The square brackets
are a tensor product of one gate acting on $A$ and one gate acting
on $B$, so they contain only single-qubit operations. The final
four-dimensional diagonal matrix cannot be
written as a product of one-qubit phase gates unless
$\Theta=0$ modulo $2\pi$. It is therefore the entangling part of
the evolution; after removing the local factors, $\Theta=\pi$
modulo $2\pi$ gives CZ.

This decomposition also identifies what must be measured. If $B$
is held in $\ket b$, the relative phase acquired by $A$ is
$d_b=\chi_{1b}-\chi_{0b}$. Hence $d_0=\alpha$,
$d_1=\alpha+\Theta$, and
\[
 \Theta=d_1-d_0.
\]

We now need an interaction that makes the logical energy of one
qubit depend on the logical state of the other. Consider two
singlet--triplet qubits $A$ and $B$, each containing two electrons,
and keep tunneling between the two qubits negligible. They still
interact electrostatically: moving charge within $A$ changes the
Coulomb energy experienced by the electrons in $B$. Because a
singlet can acquire double occupation while a triplet remains
separated, this interaction can produce the required state-dependent
energy and hence a conditional phase.

\subsubsection{Projecting the electrostatic interaction onto the logical basis}
\label{sec:calibration-capacitive-st-interaction}

We begin with one double dot before restricting it to the logical
states. Retain the charge states $S(1,1)$ and $S(0,2)$ and the
separated triplet $T_0(1,1)$. Taking the energy of $S(1,1)$ as zero
gives
\begin{equation}
 H_{\mathrm{dot}}=
 \delta_c\ket{S(0,2)}\!\bra{S(0,2)}
 +E_T\ket{T_0(1,1)}\!\bra{T_0(1,1)}
 +\sqrt2t_c
 \left[\ket{S(1,1)}\!\bra{S(0,2)}+\mathrm{h.c.}\right].
 \label{eq:calibration-st-charge-hamiltonian}
\end{equation}
The tunnel term conserves spin. It connects the two singlets, while
there is no low-energy $T_0(0,2)$ state because two electrons in the
same lowest orbital must form a spin singlet.

The two qubits have copies $H_A$ and $H_B$ of this Hamiltonian, with
their own values of $\delta_c$, $t_c$ and $E_T$. They also interact
through the ordinary Coulomb interaction. Let
$\varphi_{A,i}(\mathbf r)$ and $\varphi_{B,j}(\mathbf r)$ be the
localized orbitals on their respective dots, and let
$\hat n_{A,i}$ and $\hat n_{B,j}$ be the corresponding occupation
operators. Neglecting tunneling and orbital overlap between the two
qubits gives the density--density form
\begin{equation}
 H_C^{AB}=\sum_{i,j=1}^{2}W_{ij}\hat n_{A,i}\hat n_{B,j},
 \qquad
 W_{ij}=\int d^3r\,d^3r'\,
 |\varphi_{A,i}(\mathbf r)|^2
 V_C(\mathbf r,\mathbf r')
 |\varphi_{B,j}(\mathbf r')|^2.
 \label{eq:calibration-density-density-hamiltonian}
\end{equation}
Here $V_C$ is the electron--electron electrostatic interaction
energy, and $W_{ij}$ is the Coulomb energy of one electron in dot
$i$ of $A$ and one in dot $j$ of $B$. The Hamiltonian to be
projected is
\begin{equation}
 H=H_A+H_B+H_C^{AB}.
 \label{eq:calibration-two-st-full-hamiltonian}
\end{equation}

We next choose the logical basis of one double dot. Diagonalizing the singlet block of
Eq.~\eqref{eq:calibration-st-charge-hamiltonian} gives its lower
eigenstate
\begin{equation}
 \begin{aligned}
 \ket0=\ket{\widetilde S}
 &=\sqrt{1-p}\ket{S(1,1)}
   -\sqrt p\ket{S(0,2)},\\
 p&=\frac12\left(1-\frac{\delta_c}
                 {\sqrt{\delta_c^2+8t_c^2}}\right),
 \end{aligned}
 \qquad
 \ket1=\ket{\widetilde T_0}
 \simeq\ket{T_0(1,1)}.
 \label{eq:calibration-st-logical-charge-basis}
\end{equation}
The singlet has mean dot occupations $(1-p,1+p)$, whereas the
triplet has $(1,1)$. Thus the logical states differ mainly in spin
at the separated setting $p\simeq0$, but acquire different charge
distributions when tunneling makes $p$ nonzero. This difference
comes from spin-conserving tunneling together with the Pauli
restriction, not from the electrostatic interaction itself.

Within the logical space of one double dot, the two occupation
operators are
\begin{equation}
 \begin{aligned}
 \hat n_1
 &=(1-p)\ket0\!\bra0+\ket1\!\bra1
  =\left(1-\frac p2\right)I-\frac p2Z,\\
 \hat n_2
 &=(1+p)\ket0\!\bra0+\ket1\!\bra1
  =\left(1+\frac p2\right)I+\frac p2Z,
 \end{aligned}
 \label{eq:calibration-projected-capacitive-hamiltonian}
\end{equation}
Here $Z=\ket0\!\bra0-\ket1\!\bra1$. For qubits $A$ and $B$, use
the corresponding expressions with $(p,Z)$ replaced by
$(p_A,Z_A)$ and $(p_B,Z_B)$.

Before collecting the coefficients, separate a conditional energy
shift from those that act on only one logical qubit. Let
$E_{00},E_{01},E_{10},E_{11}$ denote the diagonal energies of the
projected Coulomb interaction in the ordered logical basis
$(\ket{00},\ket{01},\ket{10},\ket{11})$, where $\ket0$ is the
charge-hybridized singlet and $\ket1$ the separated triplet.
Define $P_{S,q}=\ket{0_q}\bra{0_q}=(I+Z_q)/2$ for $q=A,B$.
Directly matching the four diagonal entries gives
\begin{equation}
 \begin{aligned}
 \operatorname{diag}(E_{00},E_{01},E_{10},E_{11})
   ={}&E_{11}I+(E_{01}-E_{11})P_{S,A}\\
     &+(E_{10}-E_{11})P_{S,B}
       +\kappa P_{S,A}P_{S,B},\\
 \kappa={}&E_{00}-E_{01}-E_{10}+E_{11}.
 \end{aligned}
 \label{eq:st-interaction-decomposition}
\end{equation}
The first term shifts all states equally; the next two shift the
singlet energy in only one qubit. They therefore contribute global
and local phases. Expanding
$P_{S,A}P_{S,B}=(I+Z_A+Z_B+Z_AZ_B)/4$ shows that the conditional
energy supplies a $Z_AZ_B$ coefficient $\kappa/4$ as well as
additional global and local shifts.

Substituting the projected occupation operators into
Eqs.~\eqref{eq:calibration-density-density-hamiltonian} and
\eqref{eq:calibration-two-st-full-hamiltonian}, and collecting the
$I$, $Z_A$, $Z_B$ and $Z_AZ_B$ terms, gives
\begin{equation}
 \begin{aligned}
 H_{\mathrm{eff}}
 &=C I+\frac{h_A}{2}Z_A+\frac{h_B}{2}Z_B
   +\frac{\kappa}{4}Z_AZ_B,\\
 \kappa
 &=p_Ap_B(W_{11}-W_{12}-W_{21}+W_{22}).
 \end{aligned}
 \label{eq:calibration-st-capacitive-logical-hamiltonian}
\end{equation}
Here $C$ contains all state-independent energies, while $h_A$ and
$h_B$ contain the isolated-qubit splittings and their Coulomb
shifts. The result treats the interqubit Coulomb interaction as too
weak to substantially change the separate-qubit charge
eigenstates~\cite{Shulman2012}. The Hamiltonian remains diagonal while
magnetic-gradient or other terms that mix $S$ and $T_0$ have
negligible or refocused net action during the phase pulse; restoring
such terms would add $X_A$ or $X_B$ components to
$H_{\mathrm{eff}}$.

The voltage pulse controls $p_Ap_B$ and therefore $\kappa(t)$.
Start at the separated setting, where $p_A\simeq p_B\simeq0$ and
$\kappa\simeq0$. Ramp towards the two singlet charge
anticrossings, hold with $p_Ap_B>0$, and reverse the ramps. If each
branch follows its instantaneous charge eigenstate and returns to
the same logical basis state, Eq.~\eqref{eq:calibration-st-capacitive-logical-hamiltonian}
gives the conditional phase
\begin{equation}
 \Theta=-\frac1\hbar\int\kappa(t)\,dt .
 \label{eq:calibration-conditional-phase}
\end{equation}
A CZ gate requires $\Theta=\pi$ modulo $2\pi$. For a constant hold,
$|\kappa|\tau/\hbar=\pi$ is an initial duration estimate; the ramps
also contribute phase.

\subsubsection{Measuring the conditional phase and obtaining CZ}
\label{sec:calibration-conditional-phase-measurement}

The pulse of the preceding subsection leaves each logical basis
state in place but gives it a phase $\chi_{ab}$. A basis-state
population measurement cannot reveal these phases. A superposition
can: fix $B$ in $\ket b$ and prepare $A$ in
$\ket+=(\ket0+\ket1)/\sqrt2$. The full two-qubit state then evolves as
\begin{equation}
 \begin{aligned}
 \ket{\Psi_b^{\mathrm{in}}}
 &=\ket{+}_A\ket b_B,\\
 \ket{\Psi_b^{\mathrm{out}}}
 &=\frac{e^{i\chi_{0b}}}{\sqrt2}
   \bigl(\ket0_A+e^{id_b}\ket1_A\bigr)\ket b_B,
 \qquad d_b=\chi_{1b}-\chi_{0b}.
 \end{aligned}
 \label{eq:calibration-conditional-ramsey-state}
\end{equation}
Because the pulse is diagonal, $B$ remains in $\ket b$. The relative
phase between the $\ket0_A\ket b_B$ and $\ket1_A\ket b_B$ amplitudes
is $d_b$. The probabilities of the two logical-basis outcomes on $A$
remain $1/2$, but the probability of the $+1$ outcome when measuring
$X_A$ is $(1+\cos d_b)/2$. Fixing $B$ first in $\ket0$ and then
in $\ket1$ gives the
relative phases $d_0$ and $d_1$. Their difference is exactly
the phase that remains after the single-qubit factors in
Eq.~\eqref{eq:calibration-phase-factorization} are removed:
\begin{equation}
 \Theta=d_1-d_0
 =\chi_{11}-\chi_{10}-\chi_{01}+\chi_{00}.
 \label{eq:calibration-conditional-ramsey}
\end{equation}
In particular, a phase common to both choices of $b$ is a local
phase on $A$ and does not contribute to $\Theta$.

To determine each $d_b$, measure the interference in two complementary
bases. The output in Eq.~\eqref{eq:calibration-conditional-ramsey-state}
has $\langle X_A\rangle_b=\cos d_b$ and
$\langle Y_A\rangle_b=\sin d_b$. Logical-basis readout measures
$Z_A$; preceding it with $R_y(-\pi/2)$ gives the $X_A$ component,
and preceding it with $R_x(\pi/2)$ gives the $Y_A$ component.
Consequently,
\[
 d_b=\operatorname{atan2}\!\left(
   \langle Y_A\rangle_b,\langle X_A\rangle_b\right)
 \quad (\bmod\ 2\pi).
\]
Repeat this measurement with $B$ in $\ket0$ and $\ket1$ while
varying the complete pulse duration. Following the phase continuously
resolves its $2\pi$ branches. The duration is chosen where the
measured difference $\Theta$ reaches $\pi$ modulo $2\pi$; this
includes the phase accumulated during the ramps as well as the hold.

At that duration, $d_0=\alpha$ gives the phase of the local factor
on $A$. To find the local factor on $B$, instead prepare $B$ in
$\ket+$ and fix $A$ in $\ket0$. Its measured relative phase is
$\beta=\chi_{01}-\chi_{00}$. With the convention
$R_z(\phi)=e^{-i\phi Z/2}$, applying
$R_{z,A}(-\alpha)R_{z,B}(-\beta)$ after the interaction cancels
these two factors. The remaining gate is CZ, up to a global phase,
when $\Theta=\pi$ modulo $2\pi$. This inference assumes the pulse
is diagonal in the logical basis, so basis-state population transfer
must also be checked.

\subsubsection{Canceling local phases with an echoed entangling pulse}
\label{sec:calibration-echo-control}

A controlled-phase pulse generally rotates each qubit separately as
well as entangling them.
Section~\ref{sec:calibration-conditional-phase-measurement}
removes the measured local phases with subsequent rotations.
If those phases also have nearly constant but unknown offsets during
each trial, a fixed correction cannot remove them in every trial.
We instead reverse their accumulation halfway through the interaction,
without reversing the desired two-qubit term.

For an idle spin, the echo sequence of
Section~\ref{sec:measurement-ramsey-echo} cancels a different error.
A central $\pi$ pulse reverses the sign of $Z$, so a nearly constant
unwanted Zeeman offset contributes opposite phases in the two
halves. The cancellation also removes a desired phase generated
by that same $Z$ term. An idle-memory echo therefore cannot simply
be inserted into a $z$ gate without recalculating the intended
operation.

For a two-qubit diagonal interaction, echo can instead remove
local phases while retaining entanglement. Suppose the Hamiltonian
during the interaction is
\[
 H=\frac{a}{2}Z_A+\frac{b}{2}Z_B+\frac{\kappa}{4}Z_AZ_B,
\]
where $a,b,\kappa$ are constant energies during the sequence.
Simultaneous $\pi$ rotations about $x$ act, up to global phase,
as $P=X_AX_B$. Conjugation changes $Z_A$ and $Z_B$ to their
negatives but leaves $Z_AZ_B$ unchanged. Thus
\begin{align}
 PHP&=-\frac{a}{2}Z_A-\frac{b}{2}Z_B
                      +\frac{\kappa}{4}Z_AZ_B,\notag\\
 U_{\mathrm{echo}}
 &=P e^{-iH\tau/(2\hbar)}P e^{-iH\tau/(2\hbar)}
   =e^{-i\kappa\tau Z_AZ_B/(4\hbar)}.
 \label{eq:calibration-echoed-entangler}
\end{align}
The last equality follows because all three $Z$ products commute.
The two $\pi$ rotations on each qubit occur at the middle and end,
so the logical basis is restored. Local quasistatic fluctuations in
$a,b$ cancel in the same way, whereas a fluctuation in $\kappa$
does not. This construction assumes that the $\pi$ pulses are
short compared with the interaction and error time scales;
their actual durations and errors enter the subsequent
conditional-phase calibration.

To use the sequence, split the intended interaction time into two
equal halves, apply simultaneous $x$ flips at the middle and end,
and choose the total time from the measured $\kappa$.
The retained conditional phase is still $-\kappa\tau/\hbar$.
The conditional-phase measurement then tests the complete sequence,
including its finite-duration flips. Echo has canceled the local
terms; it has not corrected an error in the coupling $\kappa$.

\subsubsection{Parallel single-qubit gates in the presence of capacitive coupling}
\label{sec:calibration-parallel-st-gates}

The same electrostatic interaction matters when the intended operation
is two separate one-qubit gates. An exchange pulse on one
singlet--triplet qubit changes the charge admixture of its singlet
and hence its logical energy splitting. If both qubits receive such
pulses at once, both singlets can have different charge distributions
from their respective triplets. The Coulomb energy then depends on
the logical states of \emph{both} qubits. Applying the voltage pulses
simultaneously does not, by itself, make their joint operation a
tensor product of the two separately calibrated gates.

For exchange-dominated pulses whose charge branches return to their
logical states and whose $S$--$T_0$ mixing is negligible throughout,
the instantaneous logical Hamiltonian is the diagonal Hamiltonian
of Eq.~\eqref{eq:calibration-st-capacitive-logical-hamiltonian}.
Its coefficients can vary throughout both pulse edges and holds.
All its terms commute at different times, so, up to a global phase,
\begin{equation}
 \begin{aligned}
 U_{\parallel}
 &=R_{z,A}(\phi_A)R_{z,B}(\phi_B)
   \exp\!\left(-\frac{i\Phi_{AB}}4 Z_AZ_B\right),\\
 \phi_q&=\frac1\hbar\int h_q(t)\,dt,\qquad
 \Phi_{AB}=\frac1\hbar\int\kappa(t)\,dt=-\Theta.
 \end{aligned}
 \label{eq:calibration-parallel-st-evolution}
\end{equation}
The $h_q$ include the Coulomb-induced local shifts, so even the
one-qubit angles must be calibrated for the simultaneous waveform.
The last factor is an additional two-qubit operation. It has no
entangling part when $\Phi_{AB}=0$ modulo $2\pi$; otherwise local
$z$ corrections alone cannot turn $U_{\parallel}$ into the desired
product of one-qubit gates. This statement concerns the complete
pulse area, not just the value of $\kappa$ at the hold point. For
gates with a finite error allowance, the phase need only be close
enough to a multiple of $2\pi$.

There is a control tradeoff, but no rule that useful exchange
requires a large $\kappa$. In the charge model of
Eq.~\eqref{eq:calibration-st-charge-hamiltonian}, take
$\delta_c\gg t_c>0$ and neglect any other contribution to the
singlet--triplet splitting. The lower singlet energy shift and
double-occupation probability then give
\begin{equation}
 J_{\mathrm{tun}}\simeq\frac{2t_c^2}{\delta_c},\qquad
 p\simeq\frac{2t_c^2}{\delta_c^2}
   =\frac{J_{\mathrm{tun}}}{\delta_c}.
 \label{eq:calibration-parallel-st-tradeoff}
\end{equation}
At fixed $t_c$, moving farther from the charge anticrossing
reduces both $p$ and the exchange generated by virtual tunneling.
If the barrier and detuning can be changed together, the same
$J_{\mathrm{tun}}$ can instead be obtained at a larger $\delta_c$
and a smaller $p$, until the accessible tunnel coupling or other
device constraints limit that choice. Since this model gives
$\kappa=p_Ap_B(W_{11}-W_{12}-W_{21}+W_{22})$, specified values of
$J_A$ and $J_B$ alone do not fix the unwanted conditional phase.
Parking $B$ with $p_B\simeq0$ suppresses $\kappa$ while $A$ is
pulsed, but also reduces the tunneling-generated exchange of $B$.

For a proposed pair of simultaneous $z$ gates, the relevant question
is therefore how much $\Theta$ their actual joint waveform adds.
The conditional-phase comparison of
Section~\ref{sec:calibration-conditional-phase-measurement} can be
applied to that waveform; one can then choose operating points with
smaller charge admixture, compensate the two-qubit evolution, or
schedule the gates separately if the extra phase is too large.
When transverse $X_A$ or $X_B$ terms act appreciably during a pulse,
they do not commute with $Z_AZ_B$, and the simple phase integral in
Eq.~\eqref{eq:calibration-parallel-st-evolution} must be replaced by
the full two-qubit evolution. Experiments on GaAs
singlet--triplet arrays have measured state-dependent capacitive
shifts during simultaneous operation and have also demonstrated
parallel exchange oscillations~\cite{Yun2023,Fedele2021}. Their
coupling strengths are device dependent and do not set a value for
a silicon array.
